\begin{center}
\SourceLeaf{022b}\pdfbookmark[2]{4.}{article-004}4.
\end{center}
Pro distantia $x$ iam angulum $v$ introducamus, qui inter lineam apsidum et rectam a Sole ad corporis locum ductam (\emph{radium vectorem}) continetur, et quidem hic angulus ab ea lineae apsidum parte ubi distantiae $x$ sunt positivae incipiat, versusque eam regionem, quorsum motus corporis dirigitur, crescere supponatur. Hoc modo fit $x=r\cos v$, adeoque formula nostra $r=\frac{p}{1+e\cos v}$, unde protinus derivantur conclusiones sequentes:

I. Pro $v=0$ valor radii vectoris $r$ fit minimum, puta $=\frac{p}{1+e}$: hoc punctum \emph{perihelium} dicitur.

II. Valoribus oppositis ipsius $v$ respondent valores aequales ipsius $r$; quocirca linea apsidum sectionem conicam in duas partes aequales dirimit.

III. In \emph{ellipsi} $r$ inde a $v=0$ continuo crescit, donec valorem maximum $\frac{p}{1-e}$ assequatur in \emph{aphelio} pro $v=180^{0}$; post aphelium eodem modo rursus decrescit, quo ante increverat, donec pro $v=360^{0}$ perihelium denuo attigerit. Lineae apsidum pars perihelio hinc aphelio illinc terminata \emph{axis maior} dicitur; hinc semiaxis maior, qui etiam \emph{distantia media} vocatur, fit $=\frac{p}{1-ee}$; distantia puncti in medio axe iacentis (\emph{centri ellipsis}) a foco erit $\frac{ep}{1-ee}=ea$, denotando per $a$ semiaxem maiorem.

IV. Contra in \emph{parabola} proprie non datur aphelium, sed $r$ ultra omnes limites augetur, quo propius $v$ ad $+180^{0}$ vel $-180^{0}$ accedit. Pro $v=\pm180^{0}$ valor ipsius $r$ fit infinitus, quod indicat, curvam a linea apsidum a parte perihelio opposita non secari. Quare proprie quidem loquendo de axi maiore vel centro curvae sermo esse nequit, sed secundum analyseos usum consuetum per ampliationem formularum in ellipsi inventarum axi maiori valor infinitus tribuitur, centrumque curvae in distantia infinita a foco collocatur.

\SourceLeaf{023}V. In \emph{hyperbola} denique $v$ inter limites adhuc arctiores coërcetur, scilicet inter $v = -(180^0-\psi)$ et $v = +(180^0-\psi)$, denotando per $\psi$ angulum, cuius cosinus $= \frac{1}{e}\cdot$ Dum enim $v$ ad hosce limites appropinquat, $r$ in infinitum crescit; si vero pro $v$ alter horum limitum ipse acciperetur, valor ipsius $r$ infinitus prodiret, quod indicat, hyperbolam a recta ad lineam apsidum angulo $180^0-\psi$ supra vel infra inclinata omnino non secari. Pro valoribus hoc modo exclusis, puta a $180^0-\psi$ usque ad $180^0+\psi$, formula nostra ipsi $r$ valorem negativum assignat; recta scilicet sub tali angulo contra lineam apsidum inclinata ipsa quidem hyperbolam non secat, si vero retro producitur, in alteram hyperbolae partem incidit, quam a prima parte omnino separatam versusque eum focum quem Sol occupat convexam esse constat. Sed in disquisitione nostra, quae ut iam monuimus suppositioni innititur, $r$ sumi positive, ad hanc alteram hyperbolae partem non respiciemus, in qua corpus coeleste tale tantummodo incedere posset, in quod Sol vim non attractivam sed secundum easdem leges repulsivam exerceret. — Proprie itaque loquendo etiam in hyperbola non datur aphelium; pro aphelii analogo id partis aversae punctum, quod in linea apsidum iacet et quod respondet valoribus $v = 180^0$, $r = -\frac{p}{e-1}$, haberi poterit. Quodsi ad instar ellipsis valorem expressionis $\frac{p}{1-ee}$ etiam hic, ubi negativus evadit, semiaxem maiorem hyperbolae dicere lubet, duplum huius quantitatis puncti modo commemorati distantiam a perihelio simulque situm, ei qui in ellipsi locum habet oppositum, indicat. Perinde $\frac{ep}{1-ee}$, i. e. distantia puncti inter haec duo puncta medii (centri hyperbolae) a foco, hic obtinet valorem negativum propter situm oppositum.
